
\documentclass[12pt,a4paper]{article}

% ── Packages ─────────────────────────────────────────────────────────────
\usepackage[T1]{fontenc}
\usepackage[utf8]{inputenc}
\usepackage[margin=2.5cm]{geometry}
\usepackage{xcolor}
\usepackage{graphicx}
\usepackage{listings}
\usepackage{enumitem}
\usepackage{hyperref}
\usepackage{fancyhdr}
\usepackage{tcolorbox}
\usepackage{booktabs}
\usepackage{array}
\usepackage{tabularx}
\usepackage{titlesec}
\usepackage{mdframed}
\usepackage{fontawesome5}
\usepackage{lmodern}
\usepackage{microtype}
\usepackage{parskip}
\usepackage{longtable}

% ── Colour Palette ─────────────────────────────────────────────────────
\definecolor{mentorPurple}{HTML}{5E2D91}
\definecolor{mentorAccent}{HTML}{7B4FBC}
\definecolor{mentorGreen}{HTML}{06A77D}
\definecolor{eyantraWarn}{HTML}{F4A261}
\definecolor{eyantraDanger}{HTML}{E63946}
\definecolor{codebg}{HTML}{F4F6F9}
\definecolor{codeborder}{HTML}{CDD4E0}
\definecolor{lightgray}{HTML}{F8F9FA}
\definecolor{darkgray}{HTML}{495057}

% ── Hyperlinks ─────────────────────────────────────────────────────────
\hypersetup{
  colorlinks=true,
  linkcolor=mentorAccent,
  urlcolor=mentorAccent,
  pdftitle={EEP Auto-Grading Platform -- Mentor Guide},
  pdfauthor={E-Yantra, IIT Bombay},
}

% ── Code Listing Style ─────────────────────────────────────────────────
\lstdefinestyle{shellstyle}{
  backgroundcolor=\color{codebg},
  basicstyle=\ttfamily\small,
  breaklines=true,
  frame=single,
  rulecolor=\color{codeborder},
  framesep=4pt,
  xleftmargin=6pt,
  xrightmargin=6pt,
  showstringspaces=false,
  keywordstyle=\color{mentorPurple}\bfseries,
  commentstyle=\color{darkgray}\itshape,
}
\lstset{style=shellstyle}

% ── Custom tcolorbox environments ──────────────────────────────────────
\tcbuselibrary{skins,breakable}

\newtcolorbox{infobox}[1][Note]{
  colback=mentorAccent!8,
  colframe=mentorAccent,
  fonttitle=\bfseries,
  title={\faInfoCircle\enspace #1},
  breakable,
}

\newtcolorbox{warnbox}[1][Warning]{
  colback=eyantraWarn!12,
  colframe=eyantraWarn!80!black,
  fonttitle=\bfseries,
  title={\faExclamationTriangle\enspace #1},
  breakable,
}

\newtcolorbox{tipbox}[1][Tip]{
  colback=mentorGreen!10,
  colframe=mentorGreen,
  fonttitle=\bfseries,
  title={\faLightbulb\enspace #1},
  breakable,
}

\newtcolorbox{dangerbox}[1][Important]{
  colback=eyantraDanger!8,
  colframe=eyantraDanger,
  fonttitle=\bfseries,
  title={\faBan\enspace #1},
  breakable,
}

% ── Section Title Styling ──────────────────────────────────────────────
\titleformat{\section}
  {\Large\bfseries\color{mentorPurple}}
  {\thesection}{1em}{}[\color{mentorAccent}\rule{\linewidth}{1.5pt}]

\titleformat{\subsection}
  {\large\bfseries\color{mentorPurple}}
  {\thesubsection}{0.8em}{}

\titleformat{\subsubsection}
  {\normalsize\bfseries\color{darkgray}}
  {\thesubsubsection}{0.7em}{}

% ── Header / Footer ───────────────────────────────────────────────────
\pagestyle{fancy}
\fancyhf{}
\fancyhead[L]{\color{mentorPurple}\small\bfseries EEP Auto-Grading Platform}
\fancyhead[R]{\color{darkgray}\small Mentor Guide}
\fancyfoot[C]{\color{darkgray}\small\thepage}
\renewcommand{\headrulewidth}{0.5pt}
\renewcommand{\headrule}{\color{mentorAccent}\hrule width\headwidth height\headrulewidth}

% ── Step counter ──────────────────────────────────────────────────────
\newcounter{stepnum}
\newcommand{\step}[1]{%
  \stepcounter{stepnum}%
  \vspace{6pt}%
  \noindent\colorbox{mentorPurple}{\color{white}\bfseries\small\ Step \thestepnum\ }%
  \hspace{4pt}\textbf{#1}\\[4pt]%
}

% ─────────────────────────────────────────────────────────────────────
\begin{document}
% ─────────────────────────────────────────────────────────────────────

% ── Title Page ────────────────────────────────────────────────────────
\begin{titlepage}
  \centering
  \vspace*{2cm}

  {\color{mentorPurple}\rule{\linewidth}{3pt}}
  \vspace{0.5cm}

  {\Huge\bfseries\color{mentorPurple} EEP Auto-Grading Platform}\\[0.6em]
  {\LARGE\color{mentorAccent} Mentor Tutorial Guide}\\[0.4em]
  {\large\color{darkgray} E-Yantra Summer Internship Program (EYSIP)}

  \vspace{0.5cm}
  {\color{mentorPurple}\rule{\linewidth}{3pt}}

  \vspace{1.5cm}

  \begin{tcolorbox}[colback=lightgray, colframe=mentorAccent, width=0.85\linewidth]
    \centering
    \large This guide covers everything a mentor needs to manage\\
    \textbf{classrooms, assignments, students, and analytics}\\
    on the EEP Auto-Grading Platform.
  \end{tcolorbox}

  \vfill

  \begin{tabular}{rl}
    \textbf{Document Version:} & 1.0 \\
    \textbf{Platform Version:} & v2.0 \\
    \textbf{Audience:}         & Mentors \& Teaching Assistants \\
    \textbf{Organisation:}     & E-Yantra, IIT Bombay \\
  \end{tabular}

  \vspace{1cm}
  {\color{darkgray}\small Last updated: \today}
\end{titlepage}

% ── Table of Contents ─────────────────────────────────────────────────
\tableofcontents
\newpage

% ─────────────────────────────────────────────────────────────────────
\section{Overview}
% ─────────────────────────────────────────────────────────────────────

As a \textbf{Mentor} on the EEP Auto-Grading Platform you have elevated privileges to:

\begin{itemize}[leftmargin=*]
  \item Create and manage \textbf{Classrooms} and enrol students
  \item Create, publish, and update \textbf{Assignments}
  \item Import students in bulk via \textbf{CSV upload}
  \item Monitor \textbf{submission queues} and grading sessions
  \item View detailed \textbf{analytics} and score distributions for your assignments
  \item Export \textbf{results} and download grade reports
  \item Post \textbf{Announcements} to your students
  \item Manage \textbf{evaluator sessions} for special evaluation scenarios
\end{itemize}

\begin{infobox}[Mentor vs Admin]
  Mentors can only see and modify resources they \emph{own} (their classrooms and assignments). Admins have platform-wide access. See the Admin Guide for admin-exclusive features.
\end{infobox}

% ─────────────────────────────────────────────────────────────────────
\section{Logging In}
% ─────────────────────────────────────────────────────────────────────

\setcounter{stepnum}{0}

\step{Open the Platform}
Navigate to the platform URL in your browser. Mentors use the same login page as students.

\step{Enter Your Credentials}
\begin{itemize}
  \item \textbf{Username:} Your email address (e.g., \texttt{mentor@iitb.ac.in})
  \item \textbf{Password:} Set by the admin when your account was created, or your initial password.
\end{itemize}

\step{Confirm Your Role}
After login you will be redirected to the \textbf{Mentor Dashboard} automatically. If you see a student dashboard instead, your role may not have been assigned correctly — contact the admin.

% ─────────────────────────────────────────────────────────────────────
\section{Mentor Dashboard}
% ─────────────────────────────────────────────────────────────────────

The dashboard gives you a high-level overview of your cohort:

\begin{center}
\begin{tabularx}{0.95\linewidth}{>{\bfseries}lX}
\toprule
Widget & Description \\
\midrule
Total Students    & Students who have submitted to your assignments \\
Completion Rate   & \% of submissions that reached a terminal state \\
Average Score     & Mean score across all completed submissions \\
Total Submissions & Cumulative submission count for your assignments \\
Score Distribution & Bar chart of score bands (0--20, 21--40, \ldots, 81--100) \\
Category Breakdown & Submissions by assignment category \\
\bottomrule
\end{tabularx}
\end{center}

% ─────────────────────────────────────────────────────────────────────
\section{Managing Classrooms}
% ─────────────────────────────────────────────────────────────────────

\subsection{Creating a Classroom}

\setcounter{stepnum}{0}

\step{Navigate to Classrooms}
Click \textbf{Classrooms} in the sidebar.

\step{Create New Classroom}
\begin{enumerate}
  \item Click \textbf{+ Create Classroom}.
  \item Fill in:
    \begin{itemize}
      \item \textbf{Name} -- e.g., \texttt{Batch A -- 2025}
      \item \textbf{Description} -- A short description of the cohort
    \end{itemize}
  \item Click \textbf{Create}.
\end{enumerate}

\subsection{Managing Enrollments}

Students can self-enrol by joining your classroom code, or you can import them via CSV.

\subsubsection{Approving Enrollment Requests}

\begin{enumerate}
  \item Open the classroom and go to the \textbf{Enrollments} tab.
  \item Pending requests are shown with an \textbf{Approve} / \textbf{Reject} button.
  \item Approved students gain full access to published assignments.
\end{enumerate}

\begin{warnbox}
  Students with \emph{pending} enrollment cannot submit assignments. Approve enrollments promptly at the start of the program.
\end{warnbox}

% ─────────────────────────────────────────────────────────────────────
\section{Importing Students via CSV}
% ─────────────────────────────────────────────────────────────────────

The most efficient way to onboard a batch of students is via CSV import.

\subsection{Prepare the CSV File}

The CSV must contain exactly these three columns (order does not matter, but headers are required):

\begin{lstlisting}
full_name,email,roll_number
Jane Doe,jane@example.com,2025CS001
John Smith,john@example.com,2025CS002
\end{lstlisting}

\begin{itemize}
  \item \textbf{roll\_number} -- Must be unique. Converted to uppercase automatically.
  \item \textbf{email} -- Must be unique. Used as the student's default password.
  \item \textbf{full\_name} -- Displayed on all dashboards and reports.
  \item The file must be \textbf{UTF-8 encoded}.
\end{itemize}

\subsection{Upload the CSV}

\setcounter{stepnum}{0}

\step{Go to your Classroom}
Open the classroom you want to import students into.

\step{Import Students}
\begin{enumerate}
  \item Click \textbf{Import Students (CSV)}.
  \item Select your \texttt{.csv} file.
  \item The system processes each row:
    \begin{itemize}
      \item \textbf{created} -- New account created, enrolled in classroom (APPROVED automatically).
      \item \textbf{already\_exists} -- Student already on the platform; enrollment is still added.
    \end{itemize}
  \item A summary table is shown after import.
\end{enumerate}

\begin{tipbox}[Default Password]
  Imported students' default password is their \textbf{email address}. They are forced to change it on first login. Communicate this to your students.
\end{tipbox}

% ─────────────────────────────────────────────────────────────────────
\section{Managing Assignments}
% ─────────────────────────────────────────────────────────────────────

\subsection{Creating an Assignment}

\setcounter{stepnum}{0}

\step{Navigate to Assignments}
Click \textbf{Assignments} in the sidebar and then \textbf{+ Create Assignment}.

\step{Fill in Assignment Details}

\begin{center}
\begin{tabularx}{\linewidth}{>{\bfseries}lX}
\toprule
Field & Description \\
\midrule
Title                   & Short, descriptive name (e.g., \textit{Week 1: Git \& GitHub}) \\
Slug                    & URL-safe identifier (e.g., \texttt{week1-git}). Must be globally unique. \\
Description             & Full task description shown to students. Supports plain text. \\
Category                & One of: Git, Bash, Python, Docker, Web, ML, Other \\
Max Score               & Maximum points for full completion \\
Deadline                & Submission deadline (UTC). Cannot be changed after it passes without admin. \\
Late Penalty (\%)       & Deducted from score for submissions past the deadline (e.g., 20 = 20\% off) \\
Submission Filename     & Expected filename inside ZIP (e.g., \texttt{solution.sh}) \\
Submission Instructions & Step-by-step instructions for students \\
Expected Structure      & File tree the ZIP must follow (one path per line) \\
Expected Media URL      & Optional image/diagram URL showing expected output \\
Resource Links          & URLs to learning resources (YouTube, docs, etc.) \\
\bottomrule
\end{tabularx}
\end{center}

\step{Save as Draft}
Click \textbf{Create Assignment}. The assignment is saved as \emph{unpublished} (draft). Students cannot see it yet.

\subsection{Publishing an Assignment}

Only the assignment's \emph{creator} (or an admin) can publish it.

\setcounter{stepnum}{0}

\step{Open the assignment} from your assignments list.

\step{Click Publish} -- the assignment becomes visible to all enrolled students immediately.

\begin{dangerbox}[Before Publishing]
  Double-check the deadline, max score, and expected file structure before publishing. The deadline is \textbf{locked} after it has passed and requires admin intervention to extend.
\end{dangerbox}

\subsection{Updating an Assignment}

\begin{itemize}
  \item Click \textbf{Edit} on the assignment detail page.
  \item You can update title, description, max score, and deadline (before the deadline passes).
  \item After deadline has passed, only an admin can extend the deadline.
\end{itemize}

\subsection{Unpublishing an Assignment}

If you need to pull an assignment (e.g., an error was found), click \textbf{Unpublish}. The assignment returns to draft state and students can no longer see it. \textbf{Existing submissions are preserved.}

% ─────────────────────────────────────────────────────────────────────
\section{Monitoring Submissions}
% ─────────────────────────────────────────────────────────────────────

\subsection{Submission Queue View}

Click \textbf{Submissions} in the sidebar. You will see all submissions across your assignments:

\begin{center}
\begin{tabularx}{\linewidth}{>{\bfseries}lX}
\toprule
Column & Meaning \\
\midrule
Student         & Roll number and full name \\
Assignment      & Assignment slug and title \\
Status          & Current grading status (PENDING, RUNNING, COMPLETED, etc.) \\
Source Type     & GitHub URL or ZIP upload \\
Attempt \#      & Which attempt number this is \\
Score           & Final score (once grading completes) \\
Submitted At    & Timestamp of submission \\
\bottomrule
\end{tabularx}
\end{center}

\subsection{Results View}

Click \textbf{Results} to see only completed submissions (COMPLETED, FAILED, VALIDATION\_ERROR). This is your primary grading report view.

% ─────────────────────────────────────────────────────────────────────
\section{Analytics}
% ─────────────────────────────────────────────────────────────────────

The \textbf{Analytics} page provides insight into your cohort's performance.

\begin{itemize}
  \item \textbf{Score Distribution} -- How many students fall into each score band
  \item \textbf{Assignments Participation} -- Submission counts per assignment slug
  \item \textbf{Category Breakdown} -- Which categories (Bash, Python, etc.) are most submitted
  \item \textbf{Top Performers} -- Students with highest average scores
\end{itemize}

\begin{tipbox}[Export Results]
  Use the \textbf{Export CSV} button on the Results page to download a spreadsheet of all scores for offline analysis or uploading to your institute's grade system.
\end{tipbox}

% ─────────────────────────────────────────────────────────────────────
\section{Posting Announcements}
% ─────────────────────────────────────────────────────────────────────

\setcounter{stepnum}{0}

\step{Navigate to Announcements} Click \textbf{Announcements} in the sidebar.

\step{Create New Announcement}
\begin{enumerate}
  \item Click \textbf{+ New Announcement}.
  \item Set:
    \begin{itemize}
      \item \textbf{Title} -- Short headline
      \item \textbf{Body} -- Full announcement text
      \item \textbf{Priority} -- Normal or Urgent
    \end{itemize}
  \item Click \textbf{Post}.
\end{enumerate}

Students will see the announcement in their Announcements feed and receive an in-platform notification.

% ─────────────────────────────────────────────────────────────────────
\section{Managing Students}
% ─────────────────────────────────────────────────────────────────────

Click \textbf{Students} in the sidebar to view all students who have submitted to your assignments.

\begin{center}
\begin{tabularx}{\linewidth}{>{\bfseries}lX}
\toprule
Column & Meaning \\
\midrule
Roll Number           & Student's unique identifier \\
Full Name             & Registered name \\
Email                 & Registered email \\
Assignments Participated & Number of distinct assignments with submissions \\
Sessions Count        & Total submission count \\
\bottomrule
\end{tabularx}
\end{center}

% ─────────────────────────────────────────────────────────────────────
\section{Evaluator Sessions}
% ─────────────────────────────────────────────────────────────────────

For assignments that use external evaluator tools (e.g., EEP binary evaluators), you can start evaluator sessions from the \textbf{Evaluators} page.

\begin{itemize}
  \item Go to \textbf{Evaluators} in the sidebar.
  \item Paste the evaluator shared key (provided by the admin).
  \item Start the evaluator session for a specific student and assignment.
\end{itemize}

\begin{infobox}
  This feature is only needed for assignments that use the \texttt{.eep1} / \texttt{.eep2} proof-based evaluation scheme. Regular code submissions do not require this.
\end{infobox}

% ─────────────────────────────────────────────────────────────────────
\section{Quick Reference}
% ─────────────────────────────────────────────────────────────────────

\begin{tcolorbox}[colback=mentorPurple!5, colframe=mentorPurple, title={\faBookmark\enspace Mentor Quick Reference}]
\begin{tabular}{ll}
\textbf{Login credential:}    & Email address \\
\textbf{CSV columns:}          & \texttt{full\_name, email, roll\_number} \\
\textbf{Default student pass:} & Student's email address \\
\textbf{Publish restriction:} & Only assignment creator can publish \\
\textbf{Deadline lock:}        & Cannot extend after passing (admin required) \\
\textbf{Slug uniqueness:}      & Assignment slugs are globally unique \\
\textbf{API Docs:}             & \texttt{/docs} (if enabled) \\
\end{tabular}
\end{tcolorbox}

% ─────────────────────────────────────────────────────────────────────
\end{document}
